\documentclass[12pt,letterpaper]{article}
\usepackage{graphicx,textcomp}
\usepackage{natbib}
\usepackage{setspace}
\usepackage{fullpage}
\usepackage{color}
\usepackage[reqno]{amsmath}
\usepackage{amsthm}
\usepackage{fancyvrb}
\usepackage{amssymb,enumerate}
\usepackage[all]{xy}
\usepackage{endnotes}
\usepackage{lscape}
\newtheorem{com}{Comment}
\usepackage{float}
\usepackage{hyperref}
\newtheorem{lem} {Lemma}
\newtheorem{prop}{Proposition}
\newtheorem{thm}{Theorem}
\newtheorem{defn}{Definition}
\newtheorem{cor}{Corollary}
\newtheorem{obs}{Observation}
\usepackage[compact]{titlesec}
\usepackage{dcolumn}
\usepackage{tikz}
\usetikzlibrary{arrows}
\usepackage{multirow}
\usepackage{xcolor}
\newcolumntype{.}{D{.}{.}{-1}}
\newcolumntype{d}[1]{D{.}{.}{#1}}
\definecolor{light-gray}{gray}{0.65}
\usepackage{url}
\usepackage{listings}
\usepackage{color}

\definecolor{codegreen}{rgb}{0,0.6,0}
\definecolor{codegray}{rgb}{0.5,0.5,0.5}
\definecolor{codepurple}{rgb}{0.58,0,0.82}
\definecolor{backcolour}{rgb}{0.95,0.95,0.92}

\lstdefinestyle{mystyle}{
	backgroundcolor=\color{backcolour},   
	commentstyle=\color{codegreen},
	keywordstyle=\color{magenta},
	numberstyle=\tiny\color{codegray},
	stringstyle=\color{codepurple},
	basicstyle=\footnotesize,
	breakatwhitespace=false,         
	breaklines=true,                 
	captionpos=b,                    
	keepspaces=true,                 
	numbers=left,                    
	numbersep=5pt,                  
	showspaces=false,                
	showstringspaces=false,
	showtabs=false,                  
	tabsize=2
}
\lstset{style=mystyle}
\newcommand{\Sref}[1]{Section~\ref{#1}}
\newtheorem{hyp}{Hypothesis}

\title{Problem Set 3}
\date{Due: March 25, 2026}
\author{Zahra Khan}


\begin{document}
	\maketitle
	\section*{Instructions}
	\begin{itemize}
	\item Please show your work! You may lose points by simply writing in the answer. If the problem requires you to execute commands in \texttt{R}, please include the code you used to get your answers. Please also include the \texttt{.R} file that contains your code. If you are not sure if work needs to be shown for a particular problem, please ask.
\item Your homework should be submitted electronically on GitHub in \texttt{.pdf} form.
\item This problem set is due before 23:59 on Wednesday March 25, 2026. No late assignments will be accepted.
	\end{itemize}

	\vspace{.25cm}
\section*{Question 1}
\vspace{.25cm}
\noindent We are interested in how governments' management of public resources impacts economic prosperity. Our data come from \href{https://www.researchgate.net/profile/Adam_Przeworski/publication/240357392_Classifying_Political_Regimes/links/0deec532194849aefa000000/Classifying-Political-Regimes.pdf}{Alvarez, Cheibub, Limongi, and Przeworski (1996)} and is labelled \texttt{gdpChange.csv} on GitHub. The dataset covers 135 countries observed between 1950 or the year of independence or the first year forwhich data on economic growth are available ("entry year"), and 1990 or the last year for which data on economic growth are available ("exit year"). The unit of analysis is a particular country during a particular year, for a total $>$ 3,500 observations. 

\begin{itemize}
	\item
	Response variable: 
	\begin{itemize}
		\item \texttt{GDPWdiff}: Difference in GDP between year $t$ and $t-1$. Possible categories include: "positive", "negative", or "no change"
	\end{itemize}
	\item
	Explanatory variables: 
	\begin{itemize}
		\item
		\texttt{REG}: 1=Democracy; 0=Non-Democracy
		\item
		\texttt{OIL}: 1=if the average ratio of fuel exports to total exports in 1984-86 exceeded 50\%; 0= otherwise
	\end{itemize}
	
\end{itemize}
\newpage
\noindent Please answer the following questions:

\begin{enumerate}
	\item Construct and interpret an unordered multinomial logit with 
	\texttt{GDPWdiff} as the output and "no change" as the reference category, including the estimated cutoff points and coefficients.

The general equation used to estimate the response variable's probability in log-odds form is: 
\[ \log\left(\frac{P_{i2})}{P_{i2})}\right) = \beta_{k0} + \beta_{k1} x_1 + \cdots + \beta_{kp} x_p \]

Prior to running the model we must first prepare the data by categorizing the GDPWdiff variable as well as facotring it. Following this we must then set "No Change" as the reference category. 


	\lstinputlisting[language=R, firstline=43,lastline=48]{PS3ZK.R}
\vspace{0.1cm}
	\lstinputlisting[language=R, firstline=51,lastline=51]{PS3ZK.R}
\vspace{0.5cm}

\begin{table}[!htbp] \centering 
	\caption{} 
	\label{} 
	\begin{tabular}{@{\extracolsep{5pt}}lcc} 
		\\[-1.8ex]\hline 
		\hline \\[-1.8ex] 
		& \multicolumn{2}{c}{\textit{Dependent variable:}} \\ 
		\cline{2-3} 
		\\[-1.8ex] & Negative & Positive \\ 
		\\[-1.8ex] & (1) & (2)\\ 
		\hline \\[-1.8ex] 
		REG & 1.379$^{*}$ & 1.769$^{**}$ \\ 
		& (0.769) & (0.767) \\ 
		& & \\ 
		OIL & 4.784 & 4.576 \\ 
		& (6.885) & (6.885) \\ 
		& & \\ 
		Constant & 3.805$^{***}$ & 4.534$^{***}$ \\ 
		& (0.271) & (0.269) \\ 
		& & \\ 
		\hline \\[-1.8ex] 
		Akaike Inf. Crit. & 4,690.770 & 4,690.770 \\ 
		\hline 
		\hline \\[-1.8ex] 
		\textit{Note:}  & \multicolumn{2}{r}{$^{*}$p$<$0.1; $^{**}$p$<$0.05; $^{***}$p$<$0.01} \\ 
	\end{tabular} 
\end{table}

\begin{description}
	\item
	
\noindent Intercept Neg: On average, holding all else constant, a non-democracy with a fuel-to-total export ratio under 50\% is associated with a 3.805 increase in the log-odds of experiencing negative GDP growth relative to no change.

\noindent Intercept Neg: On average, holding all else constant, a non-democracy with a fuel-to-total export ratio under 50\% is associated with a 4.534 increase in the log-odds of experiencing positive GDP growth relative to no change.

  
\noindent Reg Negative - On average, holding all else constant, a change from a non-democracy to a democracy is associated with a 1.379 log-odds change when going from no change to a negative change in the difference in GDP.

Reg Positive - On average, holding all else constant, a change from a non-democracy to a democracy is associated with a 1.769 log-odds change when going from no change to a positive change in the difference in GDP.

Oil Negative - On average, holding all else constant, going to a country with an average fuel-to-total export ratio greater than 50\% is associated with an increase of 4.784 in log-odds when moving from no change to a negative change in the difference in GDP.

Oil Positive - On average, holding all else constant, going to a country with an average fuel-to-total export ratio greater than 50\% is associated with a 4.576 increase in log-odds when moving from no change to a positive change in the difference in GDP.

\end{description}
	\item Construct and interpret an ordered multinomial logit with \texttt{GDPWdiff} as the outcome variable, including the estimated cutoff points and coefficients.

\end{enumerate}
\begin{table}[!htbp] \centering 
	\caption{} 
	\label{} 
	\begin{tabular}{@{\extracolsep{5pt}} ccccc} 
		\\[-1.8ex]\hline 
		\hline \\[-1.8ex] 
		& Value & Std. Error & t value & p value \\ 
		\hline \\[-1.8ex] 
		REG & $0.410$ & $0.075$ & $5.456$ & $0.00000$ \\ 
		OIL & $$-$0.179$ & $0.115$ & $$-$1.549$ & $0.121$ \\ 
		No Change\textbar Negative & $$-$5.320$ & $0.252$ & $$-$21.087$ & $0$ \\ 
		Negative\textbar Positive & $$-$0.704$ & $0.048$ & $$-$14.793$ & $0$ \\ 
		\hline \\[-1.8ex] 
	\end{tabular} 
\end{table} 
Reg: Holding other variables constant, on average, for democracies, the log odds of having a positive Difference in GDP increase by about 0.410.

OIL: Holding other variables constant, on average, for countries with the ratio of fuel exports to total exports exceeding 50\%, the log odds of having a positive Difference in GDP decreases by about 0.179.

Cutoff Points: When the variable exceeds -5.320, the probability of having either a negative or positive change in GDP  becomes greater than having no change in GDP.

Cutoff Points: When the variable exceeds -0.704, the probability of having a Positive change in GDP  becomes greater than having negative change or no change.
\newpage
\begin{table}[!htbp] \centering 
	\caption{} 
	\label{} 
	\begin{tabular}{@{\extracolsep{5pt}} cccc} 
		\\[-1.8ex]\hline 
		\hline \\[-1.8ex] 
		& OR & 2.5 \% & 97.5 \% \\ 
		\hline \\[-1.8ex] 
		REG & $1.507$ & $1.301$ & $1.747$ \\ 
		OIL & $0.836$ & $0.668$ & $1.051$ \\ 
		\hline \\[-1.8ex] 
	\end{tabular} 
\end{table}
We can also interpret the results by caluclating the odds ratio. 

Reg: Holding other variables constant, on average, for democracies, the odds of having a positive Difference in GDP increases by about 50\%.

OIL: Holding other variables constant, on average, for countries with the ratio of fuel expoets to total export exceeding 50\%, the odds of having a positive Difference in GDP decreases by about 0.84\%.
\section*{Question 2} 
\vspace{.25cm}

\noindent Consider the data set \texttt{MexicoMuniData.csv}, which includes municipal-level information from Mexico. The outcome of interest is the number of times the winning PAN presidential candidate in 2006 (\texttt{PAN.visits.06}) visited a district leading up to the 2009 federal elections, which is a count. Our main predictor of interest is whether the district was highly contested, or whether it was not (the PAN or their opponents have electoral security) in the previous federal elections during 2000 (\texttt{competitive.district}), which is binary (1=close/swing district, 0="safe seat"). We also include \texttt{marginality.06} (a measure of poverty) and \texttt{PAN.governor.06} (a dummy for whether the state has a PAN-affiliated governor) as additional control variables. 

\begin{enumerate}
	\item [(a)]
	Run a Poisson regression because the outcome is a count variable. Is there evidence that PAN presidential candidates visit swing districts more? Provide a test statistic and p-value.
	
	\[
	\lambda_i = e^{\beta_0 + \beta_1 x_{i1} + \cdots + \beta_k x_{ki}}
	\]
	
	\begin{table}[!htbp] \centering 
		\caption{} 
		\label{} 
		\begin{tabular}{@{\extracolsep{5pt}}lc} 
			\\[-1.8ex]\hline 
			\hline \\[-1.8ex] 
			& \multicolumn{1}{c}{\textit{Dependent variable:}} \\ 
			\cline{2-2} 
			\\[-1.8ex] & PAN.visits.06 \\ 
			\hline \\[-1.8ex] 
			competitive.district & $-$0.081 \\ 
			& (0.171) \\ 
			& \\ 
			marginality.06 & $-$2.080$^{***}$ \\ 
			& (0.117) \\ 
			& \\ 
			PAN.governor.06 & $-$0.312$^{*}$ \\ 
			& (0.167) \\ 
			& \\ 
			Constant & $-$3.810$^{***}$ \\ 
			& (0.222) \\ 
			& \\ 
			\hline \\[-1.8ex] 
			Observations & 2,407 \\ 
			Log Likelihood & $-$645.606 \\ 
			Akaike Inf. Crit. & 1,299.213 \\ 
			\hline 
			\hline \\[-1.8ex] 
			\textit{Note:}  & \multicolumn{1}{r}{$^{*}$p$<$0.1; $^{**}$p$<$0.05; $^{***}$p$<$0.01} \\ 
		\end{tabular} 
	\end{table}
	
	(Table attached below)
We can see that there is no evidence that PAN presidential candidates visit swing districts more. the competitive. The district variable has a z-value of -0.447 and a p-value of 0.633, which, at the alpha = 0.05 threshold, is not statistically significant. 

	
	
	\item [(b)]
	Interpret the \texttt{marginality.06} and \texttt{PAN.governor.06} coefficients.
	
	Intercept: For a district which is not competitive, average poverty level, and a non-PAN governor, decreases the amount of visits from the winning PAN presidential candidate in 2006 by a multiplicative factor of 3.810.
	
	PAN.governor.06 - On average, holding all else constant, a state having a PAN-affiliated governor decreases the amount of visits from the winning PAN presidential candidate in 2006 by a multiplicative factor of 0.312. 

Marginality - On average, holding all else constant, a one-unit increase in poverty decreases the amount of visits from the winning PAN presidential candidate in 2006 by a multiplicative factor of 0.2.080. 


		\lstinputlisting[language=R, firstline=78,lastline=83]{PS3ZK.R}
	\vspace{0.5cm}
	\item [(c)]
	Provide the estimated mean number of visits from the winning PAN presidential candidate for a hypothetical district that was competitive (\texttt{competitive.district}=1), had an average poverty level (\texttt{marginality.06} = 0), and a PAN governor (\texttt{PAN.governor.06}=1).
	
	For a competitive district, with an average poverty level and a PAN governor, the estimate mean number of visitng from the winning PAN presidential candidate is about 0. 
I tried two different approaches to estimating fitted values and found that I had the same answer for both. The equation to obtain the fitted values is:

\[
\lambda_i = e^{-3.810 - 0.081 - 0.321}
\]


		\lstinputlisting[language=R, firstline=86,lastline=95]{PS3ZK.R}
\end{enumerate}

\end{document}
